\subsubsection*{Identity of the historical noise tasks}
The dataset key \texttt{noise\_resilience} was reused for two different generators. We therefore identify the generator by cohort, rather than treating the key as a single task definition. The cohort mapping, source hashes and remaining gaps are recorded in Source Data, \texttt{release\_task\_identity/task\_identity.json}.

\paragraph{Three-class noisy-line control.}
The tracked-clean 320-fit exact-transport/backpropagation rerun contains 160 MNIST fits and 160 noisy-line fits. Its frozen synthetic-data source, commit \texttt{4f3612a5}, implements the latter as follows. Draw a class $y$ uniformly from $\{0,1,2\}$ and construct a square image $L_y$: class 0 is blank, class 1 has one uniformly selected horizontal row set to one, and class 2 has one uniformly selected vertical column set to one. Each input is
\begin{equation}
 x=\operatorname{vec}(L_y+0.2\,\varepsilon),\qquad
 \varepsilon_{ab}\mathrel{\overset{\mathrm{iid}}{\sim}}\mathcal N(0,1),
 \label{eq:legacy_noisy_lines}
\end{equation}
where $a,b$ index pixels and $\operatorname{vec}$ flattens the image. The frozen dispatcher defaults to a side length of 784, producing 614,656 input coordinates; the retained checkpoint sizes independently agree with this dimension. Images are neither clipped nor normalized. The generator creates 10,000 examples and uses the run-seeded random number generators to divide them into 8,000 training, 1,000 validation and 1,000 test examples. The newer projected-noise parameters in launch configurations were not interpreted by this dispatcher. Consequently, the noise-task rows of Supplementary Table~\ref{tab:clean_exact_bp} tests agreement between two learning implementations on a three-class synthetic control, not robustness to corrupted MNIST.

\paragraph{Projected-noise MNIST protocol.}
The intended protocol for the prospective depth-by-feedback and subsequent routing, inhibitory-dose and fixed-budget cohorts uses flattened MNIST images $x\in[0,1]^{784}$ with the original ten digit labels. Let $P\in\mathbb R^{784\times50}$ have independent standard-normal entries, with each column subsequently normalized to unit Euclidean length. For each example draw $z\sim\mathcal N(0,I_{50})$ and form
\begin{equation}
 \widetilde x=\operatorname{clip}_{[0,1]}(x+1.5Pz).
 \label{eq:projected_noise_mnist}
\end{equation}
The projection seed is 0; independent latent-noise seeds are 1, 2 and 3 for the training, validation and test splits. There is no additional normalization or gain. The intended split contains 48,000 training, 12,000 validation and 10,000 test examples. Frozen prospective configurations specify these parameters, and the 784-input model sizes and accuracies in the depth-by-feedback cohort are consistent with this protocol. However, the executed nested generator file was not archived by hash. These records establish the intended protocol, not complete recovery of its historical execution. The retained lineage for the 320 fixed-budget noise fits in Supplementary Fig.~S9 joins outcomes to the input-validity audit, but its raw resolved configurations are unavailable for rechecking and do not close this generator-provenance gap.

\paragraph{Inherited aggregates and release behavior.}
The inherited synthetic feedback and gradient panels in Supplementary Figs.~S2 and S3 lack a complete join to executed configurations and generator bytes. Their task identity remains unresolved; neither generator is assigned to those aggregates retrospectively. The software release preserves the two implementations under the explicit names \texttt{legacy\_noisy\_lines} and \texttt{projected\_noise\_mnist}. It rejects the ambiguous historical key unless a generator is chosen explicitly. The projected-noise implementation is a newly frozen reference for the documented protocol, whereas the noisy-line implementation is recovered from the executed clean source commit.
